BGG community weight is a poor predictor of community satisfaction ($r = 0.187$),
confirming that game complexity and player enjoyment are largely decoupled.
This paper tests whether TIGER's five-dimensional profiles predict BGG community ratings
more accurately, using linear regression on 18 canonical games with known BGG means.
TIGER achieves $r = 0.715$ and reduces RMSE by 28.8\% relative to weight-only,
surpassing both pre-registered thresholds ($r \geq 0.60$; $\geq 15\%$ reduction).
Ablation analysis reveals that Tension (T) contributes most ($\Delta r = 0.162$),
reflecting a selection effect: high-pressure games attract committed hobbyist communities
that rate generously, while low-tension accessible games attract casual players who rate
more modestly on BGG. Elegance (E) contributes negatively ($\Delta r = 0.102$):
the highest-rated games (Brass: Birmingham, Spirit Island) are complex and hard to teach.
I, G, and R contribute minimally in this pilot sample, suggesting their predictive
signal requires larger and more varied corpora.
These findings establish that TIGER dimensions are not merely descriptive:
they capture player-relevant variation that predicts long-term community ratings.
